\documentclass[a4paper,12pt]{article}

% แพ็กเกจสำหรับภาษาไทยและการจัดการเอกสาร
\usepackage{fontspec}
\usepackage{xunicode}
\usepackage{xltxtra}
\usepackage[margin=2.5cm]{geometry}
\usepackage{enumitem}
\usepackage{listings}
\usepackage[table]{xcolor}
\usepackage{graphicx}
\usepackage{longtable}
\usepackage{booktabs}
\usepackage{tikz}
\usetikzlibrary{shapes.geometric, arrows, positioning, fit, calc}
\usepackage{float}
\usepackage[normalem]{ulem}
\usepackage{needspace}
\usepackage{amsmath}
\usepackage{amssymb}
\usepackage{multicol}
\usepackage{array}

% ตั้งค่าฟอนต์ภาษาไทย (ต้องมีฟอนต์ TH Sarabun New)
\setmainfont{TH Sarabun New}
\setmonofont{Courier New}

% ตั้งค่าการตัดคำภาษาไทย เพื่อป้องกันตัวหนังสือตกขอบกระดาษ
\XeTeXlinebreaklocale "th"
\XeTeXlinebreakskip = 0pt plus 1pt

\linespread{1.15}

\definecolor{codegreen}{rgb}{0,0.6,0}
\definecolor{codegray}{rgb}{0.5,0.5,0.5}
\definecolor{codepurple}{rgb}{0.58,0,0.82}
\definecolor{backcolour}{rgb}{0.95,0.95,0.92}

\lstdefinestyle{mystyle}{
    backgroundcolor=\color{backcolour},
    commentstyle=\color{codegreen},
    keywordstyle=\color{blue},
    numberstyle=\tiny\color{codegray},
    stringstyle=\color{codepurple},
    basicstyle=\ttfamily\footnotesize,
    breakatwhitespace=false,
    breaklines=true,
    captionpos=b,
    keepspaces=true,
    numbers=left,
    numbersep=5pt,
    showspaces=false,
    showstringspaces=false,
    showtabs=false,
    tabsize=2,
    escapeinside={(*@}{@*)}
}
\lstset{style=mystyle, language=SQL}

\newcommand{\blank}[1]{\fbox{\rule{0pt}{1.1em}\hspace{#1}}}
\newcommand{\ans}[1]{\textbf{#1}}

\begin{document}

\begin{center}
    \textbf{\Large CPE 241 Database Systems (Module 3)}\\
    \textbf{Mock Final Exam (ข้อสอบจำลองปลายภาค)}\\
    \vspace{0.2cm}
    เวลาสอบ: 2 ชั่วโมง (13.00--15.00)\\
    คะแนนเต็ม 70 คะแนน
\end{center}
\hrule
\vspace{0.35cm}

\noindent\textbf{ประกาศและขอบเขตข้อสอบ}\\
ข้อสอบฉบับนี้ออกแบบให้ครอบคลุม Lecture ทั้งหมด โดยเน้น Lecture 9--12 ได้แก่ Query Optimization, Multiuser Environment, NoSQL/MongoDB และ Advanced Topics รวมถึง SQL, Relational Algebra (RA), Tuple Relational Calculus (TRC), Domain Relational Calculus (DRC) และ Normalization\\
\textbf{สิ่งที่อนุญาตให้นำเข้าห้องสอบ:} เครื่องเขียนและบัตรนักศึกษาเท่านั้น เว้นแต่อาจารย์ผู้คุมสอบประกาศเพิ่มเติม\\
\textbf{ลักษณะข้อสอบ:} Part 1 Multiple Choice 40 ข้อ, Part 2 Short Answer 10 ข้อ, Part 3 แสดงวิธีทำ 3 ข้อ

\vspace{0.5cm}

\section*{Part 1: Multiple Choice Questions (40 Marks)}
\textbf{คำชี้แจง:} เลือกคำตอบที่ถูกต้องที่สุด (ข้อละ 1 คะแนน)

\begin{enumerate}[label=\textbf{\arabic*.}]
    \needspace{6\baselineskip}
    \item ข้อใดอธิบายลักษณะของ Relational Calculus ได้ถูกต้องที่สุด?
    \begin{enumerate}[label=\alph*.]
        \item เป็นภาษาแบบ Procedural ที่ต้องระบุลำดับการ Join ทุกขั้น
        \item เป็นภาษาแบบ Non-procedural ที่ระบุว่าต้องการข้อมูลใดโดยไม่ระบุวิธีดึงข้อมูล
        \item ใช้ได้เฉพาะกับฐานข้อมูล NoSQL เท่านั้น
        \item ใช้เพื่อสร้าง Index โดยตรงใน DBMS
    \end{enumerate}

    \needspace{6\baselineskip}
    \item ความแตกต่างหลักระหว่าง TRC และ DRC คือข้อใด?
    \begin{enumerate}[label=\alph*.]
        \item TRC ใช้ตัวแปรแทนทั้ง tuple ส่วน DRC ใช้ตัวแปรแทนค่าใน domain หรือ column
        \item TRC เป็นภาษา SQL ส่วน DRC เป็นภาษา JavaScript
        \item TRC ใช้ได้เฉพาะกับ JOIN ส่วน DRC ใช้ได้เฉพาะกับ GROUP BY
        \item TRC เป็น physical plan ส่วน DRC เป็น logical plan
    \end{enumerate}

    \needspace{6\baselineskip}
    \item Relational Algebra จัดเป็นภาษาประเภทใด และเหตุผลคืออะไร?
    \begin{enumerate}[label=\alph*.]
        \item Non-procedural เพราะระบุเฉพาะผลลัพธ์ที่ต้องการ
        \item Procedural เพราะระบุขั้นตอนการดำเนินการกับ relation เช่น selection, projection, join
        \item Declarative เพราะไม่เกี่ยวข้องกับ execution plan
        \item Schema-less เพราะไม่ต้องรู้โครงสร้าง relation
    \end{enumerate}

    \needspace{6\baselineskip}
    \item ลำดับ Query Execution Process ที่ถูกต้องที่สุดคือข้อใด?
    \begin{enumerate}[label=\alph*.]
        \item Optimization $\rightarrow$ Parsing $\rightarrow$ Execution $\rightarrow$ Binding $\rightarrow$ Result
        \item Parsing $\rightarrow$ Binding $\rightarrow$ Optimization $\rightarrow$ Execution $\rightarrow$ Result Output
        \item Binding $\rightarrow$ Execution $\rightarrow$ Parsing $\rightarrow$ Result Output $\rightarrow$ Optimization
        \item Execution $\rightarrow$ Parsing $\rightarrow$ Optimization $\rightarrow$ Binding $\rightarrow$ Result
    \end{enumerate}

    \needspace{6\baselineskip}
    \item Query Execution Plan โดยทั่วไปไม่รวมข้อมูลประเภทใด?
    \begin{enumerate}[label=\alph*.]
        \item Access method เช่น Full Table Scan, Index Scan, Index Seek
        \item Join strategy เช่น Nested Loop Join, Hash Join, Merge Join
        \item Cost estimation ของแต่ละ operation
        \item รหัสผ่านของ user ที่รัน query
    \end{enumerate}

    \needspace{6\baselineskip}
    \item จากสมการด้านล่าง เป็นเทคนิค Query Optimization แบบใด?\\
    \[
    \pi_{\text{name},\text{grade}}(\sigma_{\text{grade}='A'}(\text{STUDENT}\bowtie\text{GRADE}))
    \Rightarrow
    \pi_{\text{name},\text{grade}}(\text{STUDENT}\bowtie\sigma_{\text{grade}='A'}(\text{GRADE}))
    \]
    \begin{enumerate}[label=\alph*.]
        \item Predicate Pushdown
        \item Projection Pushdown
        \item Materialized View Rewrite
        \item Hash Partitioning
    \end{enumerate}

    \needspace{6\baselineskip}
    \item หากตาราง \texttt{STUDENT} มี 120 columns แต่ query ต้องใช้เพียง \texttt{student\_id} และ \texttt{name} ก่อนนำไป Join เทคนิคใดควรใช้เพื่อลดข้อมูลเร็วขึ้น?
    \begin{enumerate}[label=\alph*.]
        \item Projection Pushdown
        \item Dirty Read
        \item Full Outer Join
        \item Data Lakehouse
    \end{enumerate}

    \needspace{6\baselineskip}
    \item ข้อใดอธิบายเหตุผลที่ควรใช้ \texttt{WHERE} แทน \texttt{HAVING} เมื่อเงื่อนไขนั้นไม่ต้องรอผล aggregate?
    \begin{enumerate}[label=\alph*.]
        \item \texttt{WHERE} กรองก่อน GROUP BY ทำให้จำนวนแถวที่ต้อง aggregate ลดลง
        \item \texttt{WHERE} ทำงานหลัง \texttt{HAVING} เสมอ
        \item \texttt{HAVING} ใช้กับ aggregate ไม่ได้
        \item \texttt{WHERE} ทำให้ query เป็น NoSQL โดยอัตโนมัติ
    \end{enumerate}

    \needspace{6\baselineskip}
    \item Index มี trade-off หลักอย่างไร?
    \begin{enumerate}[label=\alph*.]
        \item ช่วยให้ \texttt{SELECT} เร็วขึ้น แต่เพิ่มภาระ \texttt{INSERT}, \texttt{UPDATE}, \texttt{DELETE} และใช้พื้นที่เพิ่ม
        \item ทำให้ทุก query กลายเป็น O(1) เสมอโดยไม่มีต้นทุน
        \item ใช้แทน Primary Key ได้เสมอ
        \item ทำให้ตารางเข้าสู่ 3NF โดยอัตโนมัติ
    \end{enumerate}

    \needspace{6\baselineskip}
    \item Query ต่อไปนี้ควรได้ประโยชน์จาก Index ประเภทใดมากที่สุด?\\
    \texttt{SELECT * FROM EMPLOYEE WHERE DeptID = 5 AND Salary > 50000;}
    \begin{enumerate}[label=\alph*.]
        \item Index บน \texttt{Name} เพียงอย่างเดียว
        \item Composite index บน \texttt{(DeptID, Salary)}
        \item Full-text index บน \texttt{Address}
        \item Hash index บน \texttt{Photo}
    \end{enumerate}

    \needspace{6\baselineskip}
    \item หากต้อง Join ตาราง A 1,000 แถว, B 400,000 แถว และ C 500 แถว heuristic ที่เหมาะสมคือข้อใด?
    \begin{enumerate}[label=\alph*.]
        \item Join ตารางใหญ่ที่สุดก่อนเสมอ
        \item Join ตารางเล็กก่อน เช่น A กับ C แล้วค่อย Join B
        \item ห้ามใช้ Index กับ Join column
        \item ใช้ Cartesian Product ก่อนทุกครั้ง
    \end{enumerate}

    \needspace{6\baselineskip}
    \item ลูกค้ามีจำนวนแถวจำนวนมากและต้องกระจายข้อมูลตามช่วงตัวอักษรนามสกุล A--H, I--R, S--Z ควรใช้ partitioning แบบใดตามหลักการ?
    \begin{enumerate}[label=\alph*.]
        \item Vertical Partitioning
        \item Horizontal Partitioning หรือ Range Sharding
        \item Projection Pushdown
        \item Full-text Indexing
    \end{enumerate}

    \needspace{6\baselineskip}
    \item Vertical Partitioning หมายถึงข้อใด?
    \begin{enumerate}[label=\alph*.]
        \item แบ่งตารางตามแถว เช่น ภูมิภาคหรือช่วงตัวอักษร
        \item แบ่งตารางตาม columns เช่น personal info อยู่ server หนึ่ง และ transaction history อยู่อีก server หนึ่ง
        \item สร้างสำเนาข้อมูลทั้งตารางไปทุก node
        \item บังคับให้ทุก transaction เป็น Serializable
    \end{enumerate}

    \needspace{6\baselineskip}
    \item Replication ใน Distributed Database มีประโยชน์หลักข้อใด?
    \begin{enumerate}[label=\alph*.]
        \item ลดจำนวนตารางให้เหลือตารางเดียว
        \item เพิ่ม reliability, availability และช่วยให้ระบบทนต่อ node failure ได้ดีขึ้น
        \item ทำให้ไม่ต้องใช้ concurrency control
        \item ทำให้ข้อมูลทั้งหมดกลายเป็น schema-less
    \end{enumerate}

    \needspace{6\baselineskip}
    \item ในการโอนเงินจากบัญชี A ไป B หากระบบตัดเงินจาก A แล้วดับก่อนเพิ่มเงินให้ B คุณสมบัติ ACID ใดช่วยให้ธุรกรรมต้องสำเร็จทั้งหมดหรือยกเลิกทั้งหมด?
    \begin{enumerate}[label=\alph*.]
        \item Atomicity
        \item Consistency
        \item Isolation
        \item Durability
    \end{enumerate}

    \needspace{6\baselineskip}
    \item Dirty Read หมายถึงข้อใด?
    \begin{enumerate}[label=\alph*.]
        \item อ่านข้อมูลที่ transaction อื่นยังไม่ commit และอาจ rollback ได้
        \item อ่าน row เดิมสองครั้งแล้วค่าเปลี่ยนเพราะ transaction อื่น update แล้ว commit
        \item query รอบที่สองมี row ใหม่เพิ่มเข้ามาตามเงื่อนไขเดิม
        \item อ่านข้อมูลจาก Data Lake โดยไม่มี schema
    \end{enumerate}

    \needspace{6\baselineskip}
    \item Isolation level ใดเป็นระดับต่ำสุดที่ป้องกัน Non-repeatable Read ได้ตามตาราง ANSI SQL ใน lecture?
    \begin{enumerate}[label=\alph*.]
        \item Read Uncommitted
        \item Read Committed
        \item Repeatable Read
        \item Serializable
    \end{enumerate}

    \needspace{6\baselineskip}
    \item Phantom Read ถูกป้องกันอย่างเข้มที่สุดด้วย isolation level ใด?
    \begin{enumerate}[label=\alph*.]
        \item Read Uncommitted
        \item Read Committed
        \item Repeatable Read
        \item Serializable
    \end{enumerate}

    \needspace{6\baselineskip}
    \item เมื่อ MySQL เปิด \texttt{autocommit = 1} โดย default ผลที่เกิดขึ้นคือข้อใด?
    \begin{enumerate}[label=\alph*.]
        \item ทุก SQL statement ถูก commit ทันทีหลังรันสำเร็จ ทำให้ rollback ย้อน statement นั้นไม่ได้
        \item ทุก transaction ถูก rollback อัตโนมัติหลัง SELECT
        \item DBMS จะปิด Index ทั้งหมด
        \item ห้ามใช้ \texttt{START TRANSACTION}
    \end{enumerate}

    \needspace{6\baselineskip}
    \item คำสั่ง \texttt{SAVEPOINT} ใช้ทำอะไร?
    \begin{enumerate}[label=\alph*.]
        \item สร้าง user ใหม่ในระบบ
        \item ตั้งจุดภายใน transaction เพื่อให้ rollback กลับไปยังจุดนั้นได้
        \item บังคับให้ view เป็น materialized view
        \item สร้าง schema สำหรับ MongoDB
    \end{enumerate}

    \needspace{6\baselineskip}
    \item \texttt{SELECT ... FOR UPDATE} ภายใน transaction มีผลหลักอย่างไร?
    \begin{enumerate}[label=\alph*.]
        \item อ่านข้อมูลโดยไม่ lock row ใด ๆ
        \item lock row ที่เลือกไว้เพื่อกัน transaction อื่นมา update จนกว่า commit หรือ rollback
        \item สร้าง trigger หลังอ่านข้อมูล
        \item แปลง query ให้เป็น DRC
    \end{enumerate}

    \needspace{6\baselineskip}
    \item Trigger แบบ \texttt{BEFORE INSERT} เหมาะกับงานใดมากที่สุด?
    \begin{enumerate}[label=\alph*.]
        \item Validation หรือปรับค่า \texttt{NEW} ก่อน row ถูก insert จริง
        \item สร้าง backup ของทั้ง database หลัง commit
        \item ปิด isolation level
        \item ทำ schema-on-read
    \end{enumerate}

    \needspace{6\baselineskip}
    \item ข้อใดเป็นข้อจำกัดของ Trigger ใน MySQL ตาม lecture?
    \begin{enumerate}[label=\alph*.]
        \item MySQL รองรับเฉพาะ row-level trigger ไม่รองรับ statement-level trigger
        \item Trigger ต้องเขียนด้วย Python เท่านั้น
        \item Trigger สามารถเรียก \texttt{COMMIT} และ \texttt{ROLLBACK} ได้เสมอ
        \item ตารางหนึ่งต้องมี trigger อย่างน้อย 5 ตัว
    \end{enumerate}

    \needspace{6\baselineskip}
    \item Assertion ต่างจาก Check Constraint อย่างไรในภาพรวม?
    \begin{enumerate}[label=\alph*.]
        \item Assertion สามารถนิยามเงื่อนไขข้ามหลายตารางได้ แต่หลาย DBMS ไม่รองรับโดยตรง
        \item Check Constraint ใช้ได้เฉพาะกับ MongoDB
        \item Assertion ใช้เพื่อสร้าง user account เท่านั้น
        \item ทั้งสองอย่างคือคำสั่งเดียวกันทุกประการ
    \end{enumerate}

    \needspace{6\baselineskip}
    \item ต้องการให้ user \texttt{student} อ่านและเพิ่มข้อมูลในตาราง \texttt{mydb.employee} ได้ ควรใช้คำสั่งใด?
    \begin{enumerate}[label=\alph*.]
        \item \texttt{GRANT SELECT, INSERT ON mydb.employee}\\
        \texttt{TO 'student'@'localhost';}
        \item \texttt{REVOKE SELECT, INSERT ON mydb.employee}\\
        \texttt{FROM 'student'@'localhost';}
        \item \texttt{DROP USER 'student'@'localhost';}
        \item \texttt{ROLLBACK TO student;}
    \end{enumerate}

    \needspace{6\baselineskip}
    \item View ใน SQL คืออะไร?
    \begin{enumerate}[label=\alph*.]
        \item Virtual table จากผลลัพธ์ query ซึ่งโดยทั่วไปไม่เก็บข้อมูลเอง
        \item Physical table ที่มีข้อมูลซ้ำเต็มชุดเสมอ
        \item Index ชนิดหนึ่งที่ใช้แทน Primary Key
        \item Isolation level แบบหนึ่ง
    \end{enumerate}

    \needspace{6\baselineskip}
    \item เงื่อนไขใดทำให้ View มีโอกาสเป็น updatable view ใน MySQL มากที่สุด?
    \begin{enumerate}[label=\alph*.]
        \item View อิงจากตารางเดียว และไม่มี \texttt{DISTINCT}, \texttt{GROUP BY}, aggregate หรือ join ซับซ้อน
        \item View ต้องมี \texttt{UNION} อย่างน้อยหนึ่งครั้ง
        \item View ต้องซ่อน Primary Key ทุกครั้ง
        \item View ต้อง Join อย่างน้อยสามตาราง
    \end{enumerate}

    \needspace{6\baselineskip}
    \item NoSQL หมายถึงอะไรในบริบทของ lecture?
    \begin{enumerate}[label=\alph*.]
        \item Not Only SQL เป็นกลุ่มฐานข้อมูล non-relational ที่เหมาะกับงานบางประเภท
        \item No Security Query Language
        \item ฐานข้อมูลที่ห้ามใช้ JSON
        \item ระบบสำรองข้อมูลของ RDBMS เท่านั้น
    \end{enumerate}

    \needspace{6\baselineskip}
    \item ลำดับโครงสร้างของ MongoDB จากใหญ่ไปเล็กคือข้อใด?
    \begin{enumerate}[label=\alph*.]
        \item Database $\rightarrow$ Collection $\rightarrow$ Document $\rightarrow$ Field
        \item Field $\rightarrow$ Document $\rightarrow$ Database $\rightarrow$ Collection
        \item Table $\rightarrow$ Row $\rightarrow$ Column $\rightarrow$ Index
        \item Node $\rightarrow$ Edge $\rightarrow$ Graph $\rightarrow$ Database
    \end{enumerate}

    \needspace{6\baselineskip}
    \item MongoDB Document มีลักษณะตรงกับข้อใด?
    \begin{enumerate}[label=\alph*.]
        \item Stored in BSON, รองรับ nested object, array และ field ที่ยืดหยุ่น
        \item ต้องมี schema เดียวกันทุก document แบบเข้มงวดเหมือน table ใน SQL
        \item เก็บข้อมูลได้เฉพาะ string เท่านั้น
        \item ห้ามมีค่า null และ Date
    \end{enumerate}

    \needspace{6\baselineskip}
    \item NoSQL แบบ Key-Value เหมาะกับงานใดมากที่สุด?
    \begin{enumerate}[label=\alph*.]
        \item Caching, session storage และ real-time counters
        \item ออกแบบ ERD แบบ normalized ถึง BCNF
        \item แปลง TRC เป็น DRC
        \item สร้าง OLAP cube ที่มี schema ตายตัว
    \end{enumerate}

    \needspace{6\baselineskip}
    \item Graph Database เช่น Neo4j เหมาะกับโจทย์ใด?
    \begin{enumerate}[label=\alph*.]
        \item Social network และ recommendation ที่เน้นความสัมพันธ์ระหว่าง entities
        \item Transaction bank transfer ที่ต้อง normalized สูงเท่านั้น
        \item เก็บไฟล์ raw image โดยไม่วิเคราะห์
        \item สร้าง trigger ใน MySQL
    \end{enumerate}

    \needspace{6\baselineskip}
    \item Mongoose ใน Node.js มีบทบาทอะไรกับ MongoDB?
    \begin{enumerate}[label=\alph*.]
        \item เป็น ODM ที่ช่วย model schema, validation, CRUD และ query building
        \item เป็น isolation level ของ MySQL
        \item เป็นชนิดของ Index ใน Relational DB
        \item เป็นคำสั่ง SQL สำหรับสร้าง view
    \end{enumerate}

    \needspace{6\baselineskip}
    \item ข้อใดเป็น JSON ที่ถูกต้องตามกฎพื้นฐาน?
    \begin{enumerate}[label=\alph*.]
        \item \texttt{\{ "name": "Alice", "age": 21, "isStudent": true \}}
        \item \texttt{\{ 'name': 'Alice', age: 21, \}}
        \item \texttt{\{ name = "Alice"; age = 21 \}}
        \item \texttt{[name: Alice, age: 21]}
    \end{enumerate}

    \needspace{6\baselineskip}
    \item ข้อใดจับคู่ OLTP และ OLAP ได้ถูกต้องที่สุด?
    \begin{enumerate}[label=\alph*.]
        \item OLTP: short frequent transactions; OLAP: complex analytical read-heavy queries
        \item OLTP: historical cube only; OLAP: bank transfer real-time only
        \item OLTP: schema-on-read; OLAP: no aggregation
        \item OLTP และ OLAP เหมือนกันทุกด้าน
    \end{enumerate}

    \needspace{6\baselineskip}
    \item Data Warehouse มีลักษณะเด่นข้อใด?
    \begin{enumerate}[label=\alph*.]
        \item Structured, cleaned, historical data ผ่าน ETL และ schema-on-write เพื่อ analytical queries
        \item Raw unprocessed data ทุกชนิดโดยไม่กำหนด schema จนกว่าจะอ่าน
        \item ใช้แทน transaction database เพื่อทำ bank transfer real-time
        \item ห้ามเก็บข้อมูลย้อนหลัง
    \end{enumerate}

    \needspace{6\baselineskip}
    \item Data Lake แตกต่างจาก Data Warehouse อย่างไร?
    \begin{enumerate}[label=\alph*.]
        \item Data Lake เก็บ raw/unprocessed data และใช้ schema-on-read ส่วน Data Warehouse เก็บข้อมูลที่จัดโครงสร้างและทำความสะอาดแล้ว
        \item Data Lake ใช้ได้เฉพาะกับ OLTP เท่านั้น
        \item Data Warehouse ไม่มี schema
        \item ทั้งสองอย่างใช้กับข้อมูลชนิดเดียวกันและไม่มีความต่าง
    \end{enumerate}

    \needspace{6\baselineskip}
    \item บริษัทเก็บ log clickstream, รูปภาพ, JSON event และไฟล์ raw จำนวนมากไว้ก่อน โดยยังไม่รู้ว่าจะวิเคราะห์อะไรในอนาคต ควรเลือกอะไรเหมาะสุด?
    \begin{enumerate}[label=\alph*.]
        \item Transaction Database
        \item Data Warehouse
        \item Data Lake
        \item Updatable View
    \end{enumerate}

    \needspace{6\baselineskip}
    \item ใน Distributed Database เมื่อเกิด network failure มักต้องพิจารณา trade-off ใด?
    \begin{enumerate}[label=\alph*.]
        \item Consistency, Availability และ Partition Tolerance
        \item Projection, Selection และ Rename
        \item Atomicity, Indexing และ Mongoose
        \item JSON, BSON และ XML
    \end{enumerate}

    \needspace{6\baselineskip}
    \item Elasticity ใน Cloud Computing หมายถึงข้อใด?
    \begin{enumerate}[label=\alph*.]
        \item ความสามารถในการเพิ่มหรือลด resource ตาม demand อัตโนมัติหรือรวดเร็ว
        \item การทำให้ข้อมูลทุกตัวเป็น atomic value
        \item การซ่อน column ด้วย view
        \item การ rollback transaction ไปยัง savepoint
    \end{enumerate}
\end{enumerate}

\newpage
\section*{Part 2: Short Answer Questions (10 Marks)}
\textbf{คำชี้แจง:} ตอบสั้น กระชับ และตรงประเด็น (ข้อละ 1 คะแนน)

\begin{enumerate}[label=\textbf{\arabic*.}]
    \item อธิบายความต่างระหว่าง TRC และ DRC ในแง่ของตัวแปรที่ใช้ในนิพจน์ relational calculus
    \vspace{1.3cm}

    \item จาก query \texttt{SELECT * FROM STUDENT WHERE program='CPE';} จงบอกเหตุผล 1 ข้อว่าทำไม \texttt{SELECT *} อาจทำให้ performance แย่กว่าการเลือกเฉพาะ column ที่ใช้จริง
    \vspace{1.3cm}

    \item \texttt{Index Scan} กับ \texttt{Index Seek} ต่างกันอย่างไรในภาพรวม?
    \vspace{1.3cm}

    \item ใน transaction isolation จงยกตัวอย่างสั้น ๆ ของ Phantom Read
    \vspace{1.3cm}

    \item ถ้าต้องการ undo เฉพาะบางคำสั่งภายใน transaction โดยไม่ทิ้ง transaction ทั้งหมด ควรใช้คำสั่งใดร่วมกับ \texttt{ROLLBACK TO}?
    \vspace{1.3cm}

    \item Trigger แบบ \texttt{AFTER INSERT} เหมาะกับงานประเภท logging/auditing อย่างไร?
    \vspace{1.3cm}

    \item View ช่วยเรื่อง security ได้อย่างไร จงตอบด้วยตัวอย่างสั้น ๆ
    \vspace{1.3cm}

    \item เขียนลำดับชั้นของ MongoDB ตั้งแต่ Database ลงไปถึง Field
    \vspace{1.3cm}

    \item เลือกคำตอบระหว่าง Database, Data Warehouse, Data Lake: ระบบ dashboard ผู้บริหารที่วิเคราะห์ยอดขายย้อนหลังหลายปีจากข้อมูลที่ clean แล้วควรใช้อะไร?
    \vspace{1.3cm}

    \item ตาราง \texttt{ENROLLMENT(Sid, CourseId, Sname, CourseName, Grade)} มีคีย์ผสม \texttt{\{Sid, CourseId\}} และมี \texttt{Sid $\rightarrow$ Sname} กับ \texttt{CourseId $\rightarrow$ CourseName} ขัดกับ Normal Form ระดับใดเป็นหลัก?
    \vspace{1.3cm}
\end{enumerate}

\newpage
\section*{Part 3: Long Questions (20 Marks)}
\textbf{คำชี้แจง:} แสดงวิธีทำให้ครบถ้วน ข้อ 1--2 ให้เติมช่องว่างของ SQL, RA, TRC และ DRC ให้สมบูรณ์ ส่วนข้อ 3 ให้ทำ Dependency/Normalization

\subsection*{Question 1: Query Writing 4 Forms (7 Marks)}
กำหนด schema ของระบบร้านค้าออนไลน์ดังนี้
\begin{itemize}
    \item \texttt{CUSTOMER(Cid, Cname, City, Tier)}
    \item \texttt{ORDERS(Oid, Cid, OrderDate, Status)}
    \item \texttt{ORDER\_ITEM(Oid, Pid, Qty, UnitPrice)}
    \item \texttt{PRODUCT(Pid, Pname, Category, SupplierId)}
    \item \texttt{SUPPLIER(SupplierId, Sname, Country)}
\end{itemize}

\noindent\textbf{โจทย์:} หาชื่อและ Tier ของลูกค้าที่อยู่เมือง \texttt{'Bangkok'} และเคยซื้อสินค้า Category = \texttt{'Storage'} จาก supplier ประเทศ \texttt{'Japan'} โดย order ต้องมี Status = \texttt{'PAID'}\\
\textbf{หมายเหตุ:} ต้องคืนผลลัพธ์เป็น \texttt{Cname, Tier} เท่านั้น

\vspace{0.2cm}
\textbf{1.1 SQL (2 คะแนน)}
\begin{lstlisting}[language=SQL]
SELECT DISTINCT C.Cname, C.Tier
FROM CUSTOMER C
JOIN ORDERS O ON (*@\blank{3.0cm}@*)
JOIN ORDER_ITEM I ON (*@\blank{3.0cm}@*)
JOIN PRODUCT P ON (*@\blank{3.0cm}@*)
JOIN SUPPLIER S ON (*@\blank{3.0cm}@*)
WHERE C.City = 'Bangkok'
  AND O.Status = 'PAID'
  AND P.Category = 'Storage'
  AND (*@\blank{3.0cm}@*);
\end{lstlisting}

\vspace{0.1cm}
\textbf{1.2 Relational Algebra (2 คะแนน)}
\[
\begin{aligned}
&\pi_{\text{Cname},\text{Tier}}\Big(
  \sigma_{\blank{4.3cm}}\big(\\
&\quad \text{CUSTOMER}
  \bowtie_{\blank{1.8cm}}
  \text{ORDERS}
  \bowtie_{\blank{1.8cm}}
  \text{ORDER\_ITEM}\\
&\quad \bowtie_{\blank{1.8cm}}
  \text{PRODUCT}
  \bowtie_{\blank{1.8cm}}
  \text{SUPPLIER}
  \big)\Big)
\end{aligned}
\]

\vspace{0.1cm}
\textbf{1.3 TRC (1.5 คะแนน)}
\[
\begin{aligned}
\{\, c.\text{Cname}, c.\text{Tier} \mid\;& c \in \text{CUSTOMER} \land c.\text{City}='Bangkok' \land\\
&\exists o \in \text{ORDERS}\;
\exists i \in \text{ORDER\_ITEM}\;
\exists p \in \text{PRODUCT}\;
\exists s \in \text{SUPPLIER}\\
&(\blank{10.5cm}) \,\}
\end{aligned}
\]

\vspace{0.1cm}
\textbf{1.4 DRC (1.5 คะแนน)}
\[
\begin{aligned}
\{\, cn,tier \mid\;& \exists cid,city,oid,odate,status,pid,qty,price,pn,cat,sid,sn,country\\
&(\text{CUSTOMER}(cid,cn,city,tier) \land
\text{ORDERS}(oid,cid,odate,status) \land\\
&\text{ORDER\_ITEM}(oid,pid,qty,price) \land
\text{PRODUCT}(pid,pn,cat,sid) \land\\
&\text{SUPPLIER}(sid,sn,country) \land \blank{8.5cm}) \,\}
\end{aligned}
\]

\newpage
\subsection*{Question 2: Universal Query / Division (7 Marks)}
กำหนด schema ของระบบลงทะเบียนเรียนดังนี้
\begin{itemize}
    \item \texttt{STUDENT(Sid, Sname, Major, Year)}
    \item \texttt{COURSE(Cid, Cname, Track)}
    \item \texttt{ENROLL(Sid, Cid, Grade)}
\end{itemize}

\noindent\textbf{โจทย์:} หารหัสและชื่อนักศึกษา Major = \texttt{'CPE'} ที่ลงครบทุกวิชาใน Track = \texttt{'Database'} และทุกวิชานั้นได้ Grade อยู่ใน \texttt{\{'A', 'B+', 'B'\}} หรือดีกว่าเทียบตามโจทย์นี้\\
\textbf{หมายเหตุ:} ให้ถือว่า \texttt{'A'}, \texttt{'B+'}, \texttt{'B'} เป็นกลุ่มเกรดผ่านเงื่อนไข

\vspace{0.2cm}
\textbf{2.1 SQL แบบ \texttt{NOT EXISTS} ซ้อนกัน (2 คะแนน)}
\begin{lstlisting}[language=SQL]
SELECT S.Sid, S.Sname
FROM STUDENT S
WHERE S.Major = 'CPE'
  AND NOT EXISTS (
      SELECT 1
      FROM COURSE C
      WHERE C.Track = 'Database'
        AND NOT EXISTS (
            SELECT 1
            FROM ENROLL E
            WHERE (*@\blank{3.0cm}@*)
              AND (*@\blank{3.0cm}@*)
              AND (*@\blank{4.0cm}@*)
        )
  );
\end{lstlisting}

\vspace{0.1cm}
\textbf{2.2 Relational Algebra แบบ Division (2 คะแนน)}
\[
\begin{aligned}
DBCourse &= \pi_{\text{Cid}}(\sigma_{\text{Track}='Database'}(\text{COURSE}))\\
GoodEnroll &= \pi_{\text{Sid},\text{Cid}}(\sigma_{\blank{3.8cm}}(\text{ENROLL}))\\
GoodStudent &= \blank{5.2cm}\\
\text{Answer} &= \pi_{\text{Sid},\text{Sname}}(\sigma_{\text{Major}='CPE'}(\text{STUDENT}) \bowtie_{\blank{2.2cm}} GoodStudent)
\end{aligned}
\]

\vspace{0.1cm}
\textbf{2.3 TRC (1.5 คะแนน)}
\[
\begin{aligned}
\{\, s.\text{Sid},s.\text{Sname} \mid\;& s \in \text{STUDENT} \land s.\text{Major}='CPE' \land\\
&\forall c \in \text{COURSE}\,
(c.\text{Track}='Database' \Rightarrow\\
&\quad \exists e \in \text{ENROLL}\,(\blank{9.5cm})) \,\}
\end{aligned}
\]

\vspace{0.1cm}
\textbf{2.4 DRC (1.5 คะแนน)}
\[
\begin{aligned}
\{\, sid,sname \mid\;& \exists major,yr(
\text{STUDENT}(sid,sname,major,yr) \land major='CPE' \land\\
&\forall cid,cname,track(\text{COURSE}(cid,cname,track) \land track='Database' \Rightarrow\\
&\quad \exists grade(\text{ENROLL}(sid,cid,grade) \land \blank{4.4cm}))) \,\}
\end{aligned}
\]

\newpage
\subsection*{Question 3: Functional Dependency and Normalization (6 Marks)}
กำหนด relation สำหรับบันทึกงานโปรเจกต์ดังนี้
\begin{center}
\small\ttfamily
\begin{tabular}{l}
PROJECT\_ASSIGNMENT(EmpID, ProjectID, EmpName, DeptID, DeptName,\\
\quad ProjectName, ClientID, ClientName,\\
\quad ManagerID, ManagerName, Hours, SkillLevel)
\end{tabular}
\end{center}
กำหนด Primary Key เป็น \texttt{\{EmpID, ProjectID\}} และมี Functional Dependencies:
\begin{itemize}
    \item \textbf{FD1:} \texttt{EmpID $\rightarrow$ EmpName, DeptID}
    \item \textbf{FD2:} \texttt{DeptID $\rightarrow$ DeptName}
    \item \textbf{FD3:} \texttt{ProjectID $\rightarrow$ ProjectName, ClientID, ManagerID}
    \item \textbf{FD4:} \texttt{ClientID $\rightarrow$ ClientName}
    \item \textbf{FD5:} \texttt{ManagerID $\rightarrow$ ManagerName}
    \item \textbf{FD6:} \texttt{\{EmpID, ProjectID\} $\rightarrow$ Hours, SkillLevel}
\end{itemize}

\begin{enumerate}[label=\textbf{3.\arabic*}]
    \item (1 คะแนน) Relation เดิมอยู่ใน Normal Form สูงสุดระดับใด และขัดกับกฎใดของระดับถัดไป? ให้ระบุ dependency ที่เป็นปัญหา
    \vspace{2.5cm}

    \item (2 คะแนน) แยก relation ให้ผ่าน 2NF โดยระบุ schema, Primary Key และ Foreign Key
    \vspace{4.0cm}

    \item (3 คะแนน) แยกต่อให้ผ่าน 3NF โดยระบุ schema สุดท้ายทั้งหมด, Primary Key และ Foreign Key พร้อมอธิบายว่าแก้ transitive dependency ตัวใดบ้าง
    \vspace{5.0cm}
\end{enumerate}

\newpage
\section*{เฉลยและแนวคิด (Answer Key)}
\textit{ส่วนนี้สำหรับตรวจหลังทำข้อสอบ แนะนำให้ทำข้อสอบเต็มเวลาก่อนเปิดเฉลย}

\subsection*{Part 1: Multiple Choice Questions}
\begin{multicols}{2}
\begin{enumerate}[label=\textbf{\arabic*.}]
    \item \ans{b.} Relational Calculus เป็น Non-procedural
    \item \ans{a.} TRC ใช้ tuple variables, DRC ใช้ domain variables
    \item \ans{b.} RA เป็น Procedural
    \item \ans{b.} Parsing, Binding, Optimization, Execution, Result
    \item \ans{d.} Query plan ไม่แสดง password
    \item \ans{a.} Predicate Pushdown
    \item \ans{a.} Projection Pushdown
    \item \ans{a.} WHERE กรองก่อน aggregate
    \item \ans{a.} Index ช่วยอ่านแต่มีต้นทุนเขียนและ storage
    \item \ans{b.} Composite index บน \texttt{(DeptID, Salary)}
    \item \ans{b.} Join ตารางเล็กก่อนตาม heuristic
    \item \ans{b.} แบ่งตามแถวคือ Horizontal หรือ Range Sharding
    \item \ans{b.} Vertical คือแบ่งตาม columns
    \item \ans{b.} Replication ช่วย reliability และ fault tolerance
    \item \ans{a.} Atomicity
    \item \ans{a.} Dirty Read คืออ่านค่าที่ยังไม่ commit
    \item \ans{c.} Repeatable Read
    \item \ans{d.} Serializable
    \item \ans{a.} Autocommit commit ทุก statement
    \item \ans{b.} SAVEPOINT ตั้งจุด rollback ภายใน transaction
    \item \ans{b.} SELECT FOR UPDATE lock row
    \item \ans{a.} BEFORE INSERT ใช้ validate/modify ก่อน insert
    \item \ans{a.} MySQL มี row-level trigger
    \item \ans{a.} Assertion ใช้เงื่อนไขข้าม table ได้ แต่หลาย DBMS ไม่รองรับ
    \item \ans{a.} GRANT SELECT, INSERT
    \item \ans{a.} View เป็น virtual table
    \item \ans{a.} Updatable view มักต้องอิง table เดียวและไม่ aggregate/join ซับซ้อน
    \item \ans{a.} NoSQL = Not Only SQL
    \item \ans{a.} Database, Collection, Document, Field
    \item \ans{a.} MongoDB เก็บ BSON และ nested/array ได้
    \item \ans{a.} Key-value เหมาะกับ caching/session
    \item \ans{a.} Graph เหมาะกับ relationship-heavy use cases
    \item \ans{a.} Mongoose เป็น ODM
    \item \ans{a.} JSON ต้องใช้ double quotes และไม่มี trailing comma
    \item \ans{a.} OLTP transaction สั้นถี่, OLAP analytical read-heavy
    \item \ans{a.} Data Warehouse คือ structured historical analytical store
    \item \ans{a.} Data Lake เป็น raw/schema-on-read
    \item \ans{c.} raw diverse future analytics คือ Data Lake
    \item \ans{a.} CAP trade-off
    \item \ans{a.} Elasticity คือเพิ่มลด resource ตาม demand
\end{enumerate}
\end{multicols}

\subsection*{Part 2: Short Answer Questions}
\begin{enumerate}[label=\textbf{\arabic*.}]
    \item TRC ใช้ตัวแปรแทนทั้ง tuple/row เช่น \texttt{t.Attribute}; DRC ใช้ตัวแปรแทนค่าในแต่ละ domain/column
    \item \texttt{SELECT *} ดึง column เกินจำเป็น ทำให้ data transfer สูงขึ้น และลดโอกาสทำ projection pushdown ได้ดี
    \item \texttt{Index Scan} ไล่อ่านช่วง/ส่วนของ index หลาย entry; \texttt{Index Seek} กระโดดไปยัง entry ที่เกี่ยวข้องโดยตรงมากกว่า จึงมักมีต้นทุนต่ำกว่า
    \item ตัวอย่าง: query \texttt{SELECT * FROM orders WHERE price > 1000} สองรอบใน transaction เดียวกัน รอบสองเจอ row ใหม่ที่ transaction อื่น insert และ commit เข้ามา
    \item ใช้ \texttt{SAVEPOINT} เช่น \texttt{SAVEPOINT sp1;} แล้ว \texttt{ROLLBACK TO sp1;}
    \item ใช้ insert record ลงตาราง log หลัง row ใหม่ถูก insert สำเร็จ เช่น เก็บ \texttt{emp\_id}, action, timestamp
    \item สร้าง view ที่เลือกเฉพาะ column ปลอดภัย เช่น \texttt{id, name, position} แล้ว grant ให้ user โดยไม่ให้เห็น \texttt{salary}
    \item Database $\rightarrow$ Collection $\rightarrow$ Document $\rightarrow$ Field
    \item Data Warehouse
    \item ขัดกับ 2NF เพราะมี Partial Dependency จากส่วนหนึ่งของ composite key ไปยัง non-prime attributes
\end{enumerate}

\subsection*{Part 3: Long Questions}

\subsubsection*{Question 1}
\textbf{SQL}
\begin{lstlisting}[language=SQL]
SELECT DISTINCT C.Cname, C.Tier
FROM CUSTOMER C
JOIN ORDERS O ON C.Cid = O.Cid
JOIN ORDER_ITEM I ON O.Oid = I.Oid
JOIN PRODUCT P ON I.Pid = P.Pid
JOIN SUPPLIER S ON P.SupplierId = S.SupplierId
WHERE C.City = 'Bangkok'
  AND O.Status = 'PAID'
  AND P.Category = 'Storage'
  AND S.Country = 'Japan';
\end{lstlisting}

\textbf{RA}
\[
\begin{aligned}
&\pi_{\text{Cname},\text{Tier}}\Big(
  \sigma_{\text{City}='Bangkok' \land \text{Status}='PAID' \land \text{Category}='Storage' \land \text{Country}='Japan'}\big(\\
&\quad \text{CUSTOMER}
  \bowtie_{\text{CUSTOMER.Cid}=\text{ORDERS.Cid}}
  \text{ORDERS}
  \bowtie_{\text{ORDERS.Oid}=\text{ORDER\_ITEM.Oid}}
  \text{ORDER\_ITEM}\\
&\quad \bowtie_{\text{ORDER\_ITEM.Pid}=\text{PRODUCT.Pid}}
  \text{PRODUCT}
  \bowtie_{\text{PRODUCT.SupplierId}=\text{SUPPLIER.SupplierId}}
  \text{SUPPLIER}
  \big)\Big)
\end{aligned}
\]

\textbf{TRC}
\[
\begin{aligned}
\{\, c.\text{Cname}, c.\text{Tier} \mid\;& c \in \text{CUSTOMER} \land c.\text{City}='Bangkok' \land\\
&\exists o \in \text{ORDERS}\;
\exists i \in \text{ORDER\_ITEM}\;
\exists p \in \text{PRODUCT}\;
\exists s \in \text{SUPPLIER}\\
&(o.\text{Cid}=c.\text{Cid} \land i.\text{Oid}=o.\text{Oid} \land p.\text{Pid}=i.\text{Pid} \land\\
&\quad s.\text{SupplierId}=p.\text{SupplierId} \land o.\text{Status}='PAID' \land\\
&\quad p.\text{Category}='Storage' \land s.\text{Country}='Japan') \,\}
\end{aligned}
\]

\textbf{DRC}
\[
\begin{aligned}
\{\, cn,tier \mid\;& \exists cid,city,oid,odate,status,pid,qty,price,pn,cat,sid,sn,country\\
&(\text{CUSTOMER}(cid,cn,city,tier) \land
\text{ORDERS}(oid,cid,odate,status) \land\\
&\text{ORDER\_ITEM}(oid,pid,qty,price) \land
\text{PRODUCT}(pid,pn,cat,sid) \land\\
&\text{SUPPLIER}(sid,sn,country) \land city='Bangkok' \land status='PAID' \land\\
&cat='Storage' \land country='Japan') \,\}
\end{aligned}
\]

\subsubsection*{Question 2}
\textbf{SQL}
\begin{lstlisting}[language=SQL]
SELECT S.Sid, S.Sname
FROM STUDENT S
WHERE S.Major = 'CPE'
  AND NOT EXISTS (
      SELECT 1
      FROM COURSE C
      WHERE C.Track = 'Database'
        AND NOT EXISTS (
            SELECT 1
            FROM ENROLL E
            WHERE E.Sid = S.Sid
              AND E.Cid = C.Cid
              AND E.Grade IN ('A', 'B+', 'B')
        )
  );
\end{lstlisting}

\textbf{RA}
\[
\begin{aligned}
DBCourse &= \pi_{\text{Cid}}(\sigma_{\text{Track}='Database'}(\text{COURSE}))\\
GoodEnroll &= \pi_{\text{Sid},\text{Cid}}(\sigma_{\text{Grade}\in\{'A','B+','B'\}}(\text{ENROLL}))\\
GoodStudent &= GoodEnroll \div DBCourse\\
\text{Answer} &= \pi_{\text{Sid},\text{Sname}}(\sigma_{\text{Major}='CPE'}(\text{STUDENT}) \bowtie_{\text{STUDENT.Sid}=\text{GoodStudent.Sid}} GoodStudent)
\end{aligned}
\]

\textbf{TRC}
\[
\begin{aligned}
\{\, s.\text{Sid},s.\text{Sname} \mid\;& s \in \text{STUDENT} \land s.\text{Major}='CPE' \land\\
&\forall c \in \text{COURSE}\,
(c.\text{Track}='Database' \Rightarrow\\
&\quad \exists e \in \text{ENROLL}\,(e.\text{Sid}=s.\text{Sid} \land e.\text{Cid}=c.\text{Cid} \land e.\text{Grade}\in\{'A','B+','B'\})) \,\}
\end{aligned}
\]

\textbf{DRC}
\[
\begin{aligned}
\{\, sid,sname \mid\;& \exists major,yr(
\text{STUDENT}(sid,sname,major,yr) \land major='CPE' \land\\
&\forall cid,cname,track(\text{COURSE}(cid,cname,track) \land track='Database' \Rightarrow\\
&\quad \exists grade(\text{ENROLL}(sid,cid,grade) \land grade\in\{'A','B+','B'\}))) \,\}
\end{aligned}
\]

\subsubsection*{Question 3}
\textbf{3.1} Relation เดิมอยู่ใน \textbf{1NF} เท่านั้น เพราะมี Partial Dependency จากส่วนหนึ่งของ composite key ไปยัง non-prime attributes จึงไม่ผ่าน 2NF
\begin{itemize}
    \item Partial Dependency: \texttt{EmpID} $\rightarrow$ \texttt{EmpName, DeptID}
    \item Partial Dependency: \texttt{ProjectID} $\rightarrow$ \texttt{ProjectName, ClientID, ManagerID}
    \item Transitive Dependency: \texttt{DeptID} $\rightarrow$ \texttt{DeptName}, \texttt{ClientID} $\rightarrow$ \texttt{ClientName}, \texttt{ManagerID} $\rightarrow$ \texttt{ManagerName}
\end{itemize}

\vspace{0.2cm}
\textbf{3.2 2NF Decomposition}
\begin{itemize}
    \item \textbf{EMPLOYEE\_2NF}(\uline{EmpID}, EmpName, DeptID, DeptName)
    \item \textbf{PROJECT\_2NF}(\uline{ProjectID}, ProjectName, ClientID, ClientName, ManagerID, ManagerName)
    \item \textbf{ASSIGNMENT}(\uline{EmpID} (FK), \uline{ProjectID} (FK), Hours, SkillLevel)
\end{itemize}
เหตุผล: แยกข้อมูลที่ขึ้นกับ \texttt{EmpID} และ \texttt{ProjectID} ออกจากตารางกลาง เหลือข้อมูลที่ขึ้นกับ key ทั้งก้อนใน \texttt{ASSIGNMENT}

\vspace{0.2cm}
\textbf{3.3 3NF Decomposition}
\begin{itemize}
    \item \textbf{EMPLOYEE}(\uline{EmpID}, EmpName, DeptID (FK))
    \item \textbf{DEPARTMENT}(\uline{DeptID}, DeptName)
    \item \textbf{PROJECT}(\uline{ProjectID}, ProjectName, ClientID (FK), ManagerID (FK))
    \item \textbf{CLIENT}(\uline{ClientID}, ClientName)
    \item \textbf{MANAGER}(\uline{ManagerID}, ManagerName)
    \item \textbf{ASSIGNMENT}(\uline{EmpID} (FK), \uline{ProjectID} (FK), Hours, SkillLevel)
\end{itemize}
แก้ transitive dependencies ต่อไปนี้:
\begin{itemize}
    \item \texttt{EmpID} $\rightarrow$ \texttt{DeptID} $\rightarrow$ \texttt{DeptName}
    \item \texttt{ProjectID} $\rightarrow$ \texttt{ClientID} $\rightarrow$ \texttt{ClientName}
    \item \texttt{ProjectID} $\rightarrow$ \texttt{ManagerID} $\rightarrow$ \texttt{ManagerName}
\end{itemize}

\end{document}
